\begin{lemma}[A terminal bridge gives a finite shift]
If $\mathsf{FrontBridge}(F,s,t,u)$ and both $s,t$ belong to $F$, then
$s\mathrel{\mathsf{FiniteShift}}t$.
\end{lemma}
\begin{proof}
By the bridge clauses in \noderef{FrontBridge}, together with $hs$ and
$ht$, we have
\[
  \mathsf{InitialSegment}(s,u)
  \qquad\text{and}\qquad
  \mathsf{InitialSegment}\bigl(t,\mathsf{FrontBridgeTail}(u)\bigr).
\]
The same bridge statement gives that $u$ is nonempty.  Unpacking the two
displayed relations using \noderef{InitialSegment}, each is either equality
or the cut below a member of its right-hand finite set.

Put
\[
  X=u\cup\{k\in\mathbb N\mid\max(u)<k\}.
\]
This set is infinite, since it contains every natural number greater than
the finite bound $\max(u)$.  By \noderef{InfiniteSetEnumeration}, choose
$x\in\mathsf{IncSeq}$ with $\operatorname{range}(x)=X$.  Every member of
$u$ is at most $\max(u)$, while every member of $X\setminus u$ is greater
than $\max(u)$.  Thus $x$ lists all members of $u$ first, followed by the
continuation above $\max(u)$.

We first prove the required prefix assertion for $s$.  If $s=u$, use the
witness $\max(u)+1$.  It belongs to $X$, and for every $k\in X$,
\[
  k<\max(u)+1\quad\Longleftrightarrow\quad k\in u,
\]
because a member of $X\setminus u$ is strictly greater than $\max(u)$.
Hence $s=\{k\in\operatorname{range}(x)\mid k<\max(u)+1\}$.  If instead
\noderef{InitialSegment} gives $a\in u$ and
$s=\{k\in u\mid k<a\}$, then $a\in X$ and, for every $k\in X$,
\[
  k<a\quad\Longleftrightarrow\quad k\in u\text{ and }k<a,
\]
since $a\leq\max(u)$.  Thus in this case $a$ is a witness as well.  By
the definition of \noderef{ProperPrefixSet},
\[
  \mathsf{ProperPrefixSet}\bigl(s,\operatorname{range}(x)\bigr).
\]

Now unfold \noderef{FrontBridgeTail}.  Since $u$ is nonempty, it is exactly
$u$ with its least member removed.  By the definition of
\noderef{ShiftMap}, $\operatorname{range}(\mathsf{ShiftMap}(x))$ is the
range of the increasing enumeration after its first value is removed.
The first value is the least member of $X$, namely the least member of
$u$.  Consequently,
\[
  \operatorname{range}(\mathsf{ShiftMap}(x))
  =\mathsf{FrontBridgeTail}(u)\cup
    \{k\in\mathbb N\mid\max(u)<k\}.
\]
Use the second initial-segment relation above.  If
$t=\mathsf{FrontBridgeTail}(u)$, the witness $\max(u)+1$ belongs to the
shifted range and its elements below it are exactly
$\mathsf{FrontBridgeTail}(u)$.  Otherwise, \noderef{InitialSegment} gives
$a\in\mathsf{FrontBridgeTail}(u)$ and
$t=\{k\in\mathsf{FrontBridgeTail}(u)\mid k<a\}$.  The value $a$ belongs
to the shifted range, and no continuation value is below $a$, so the
elements of the shifted range below $a$ are exactly the displayed finite
set.  In either case the explicit witness and the definition of
\noderef{ProperPrefixSet} give
\[
  \mathsf{ProperPrefixSet}\bigl(t,
    \operatorname{range}(\mathsf{ShiftMap}(x))\bigr).
\]
The two proper-prefix assertions are precisely the conjuncts in the
definition of \noderef{FiniteShift}; therefore
$s\mathrel{\mathsf{FiniteShift}}t$.
\end{proof}
